\documentclass{article}

\usepackage[margin=1in]{geometry}
\usepackage{amsmath,amsthm,amssymb}
\usepackage[spanish]{babel}
\usepackage{enumitem}

\theoremstyle{plain}
\newtheorem{proposicion}{Proposición}[section]

\theoremstyle{definition}
\newtheorem{definition}{Definición}[section]

% Number sets
\newcommand{\R}{\mathbb{R}}
\newcommand{\Z}{\mathbb{Z}}
\newcommand{\N}{\mathbb{N}}
\newcommand{\Q}{\mathbb{Q}}
\newcommand{\C}{\mathbb{C}}

% Topology operators
\newcommand{\cl}[1]{\overline{#1}}              % Closure
\newcommand{\Int}[1]{{#1}^\circ}                % Interior
\newcommand{\bd}{\partial}                      % Boundary
\newcommand{\Ent}[1]{\mathcal{O}(#1)}           % Neighborhoods (Entornos)
\newcommand{\pf}{\mathcal{P}_F}                 % Finite parts
\newcommand{\partes}{\mathcal{P}}               % Power set

% Useful shortcuts
\newcommand{\imp}{\implies}                      % Implies

% Net convergence/accumulation notation
\newcommand{\evto}[1]{\stackrel{E}{\hookrightarrow} #1}   % Eventually in
\newcommand{\frto}[1]{\stackrel{F}{\hookrightarrow} #1}   % Frequently in
\newcommand{\nnevto}[1]{\not\!\stackrel{E}{\hookrightarrow} #1}
\newcommand{\nnfrto}[1]{\not\!\stackrel{F}{\hookrightarrow} #1}

% Exercise environment
\newenvironment{ejercicio}[1]{\begin{trivlist}
\item[\hskip \labelsep {\bfseries Ejercicio}\hskip \labelsep {\bfseries #1.}]}{\end{trivlist}}

\setlist[enumerate,1]{label=(\alph*), leftmargin=*, itemsep=2pt}

\begin{document}

\title{Topología: Práctica VIII}
\author{Bustos Jordi\\Homotopías}

\maketitle

\begin{ejercicio}{1}
    Probar que cualesquiera dos rotaciones en \( S^{1} \) son homotópicas.
\end{ejercicio}
\begin{proof}
    Sabemos que una rotación en \( S^1 \) se puede representar como \( R_\theta(z) = e^{i\theta}z \) para algún ángulo \( \theta \). Sea \( R_{\theta_1} \) y \( R_{\theta_2} \) dos rotaciones en \( S^1 \). Definimos una homotopía \( H: S^1 \times [0,1] \to S^1 \) por:
    \[
        H(z, t) = e^{i((1-t)\theta_1 + t\theta_2)}z
    \]
    Para \( t=0 \), tenemos \( H(z, 0) = e^{i\theta_1}z = R_{\theta_1}(z) \), y para \( t=1 \), tenemos \( H(z, 1) = e^{i\theta_2}z = R_{\theta_2}(z) \). Además, \( H \) es continua en ambos argumentos, por lo que \( H \) es una homotopía entre \( R_{\theta_1} \) y \( R_{\theta_2} \). Por lo tanto, cualquier par de rotaciones en \( S^1 \) son homotópicas.
\end{proof}

\vspace{0.5cm}

\begin{definition}
    \( X \) es contractible si \( id_X \) es homotópica a una función constante.
\end{definition}

\begin{ejercicio}{2}
    Probar que toda \( f : X \rightarrow S^{2} \) continua que no es suryectiva es homotópicamente nula.
\end{ejercicio}
\begin{proof}
    Como \( f \) no es sobreyectiva, existe algún punto \( p \in S^{2} \) tal que \( p \notin \text{Im}(f) \). Por lo tanto, la imagen de \( f \) está contenida en \( S^{2} \setminus \{p\} \).

    Sabemos que la esfera menos un punto, \( S^{2} \setminus \{p\} \), es homeomorfa al plano \( \R^{2} \) mediante la proyección estereográfica desde el polo \( p \). Sea \( \phi : S^{2} \setminus \{p\} \rightarrow \R^{2} \) dicho homeomorfismo.

    Podemos considerar la composición \( g = \phi \circ f \), que es una función continua de \( X \) en \( \R^{2} \). Dado que \( \R^{2} \) es un conjunto convexo (y por ende, contractible), toda función continua hacia \( \R^{2} \) es homotópicamente nula.
    En particular, podemos definir la homotopía de línea recta hacia el origen:
    \[ H : X \times {[0 , 1]} \rightarrow \R^{2} \quad \text{dada por} \quad H(x,t) = (1-t)g(x) \]
    Notemos que \( H(x,0) = g(x) \) y \( H(x,1) = 0 \) para todo \( x \in X \).

    Para llevar esta homotopía de regreso a la esfera, definimos \( F : X \times [0,1] \rightarrow S^{2} \) como:
    \[ F(x,t) = \phi^{-1}(H(x,t)) \]
    La función \( F \) es continua por ser composición de funciones continuas. Además, tenemos que:
    \[ F(x,0) = \phi^{-1}(H(x,0)) = \phi^{-1}(g(x)) = \phi^{-1}(\phi(f(x))) = f(x) \]
    \[ F(x,1) = \phi^{-1}(H(x,1)) = \phi^{-1}(0) = c \]
    donde \( c \in S^{2} \) es el punto antípoda a \( p \) (que corresponde al origen bajo \( \phi \)).

    Por lo tanto, \( F \) es una homotopía entre \( f \) y la función constante \( c \), lo que demuestra que \( f \) es homotópicamente nula.
\end{proof}

\clearpage

\begin{definition}
    Un espacio \( X \) es simplemente conexo si es conexo por caminos y el grupo fundamental \( \pi_1(X, x_0) \) es trivial para algún \( x_0 \) y, por lo tanto, para todo \( x_0 \in X \).
\end{definition}

\begin{ejercicio}{3}
    Probar que, en un espacio simplemente conexo, dos curvas cualesquiera con los mismos puntos inicial y final son \( p \)-homotópicas.
\end{ejercicio}
\begin{proof}
    Sean \( \gamma \) y \( \rho \) curvas en \( X \) simplemente conexo con \( \gamma(0) = \rho(0) \) y \( \gamma(1) = \rho(1) \).
    Por ser \( X \) simplemente conexo vale que:
    \[
        [\gamma] \ast [\overline{\rho}] = [c_{\gamma(0)}] \in \pi_1(X, \gamma(0)) \implies [\gamma] = [\rho] \in \pi_1(X, \gamma(0))
    \]
    Por lo tanto, \( \gamma \) y \( \rho \) son \( p \)-homotópicas.
\end{proof}

\vspace{0.5cm}

\begin{definition}
    Decimos que \( p : E \to B \) es un revestimiento si \( p \) es continua, suryectiva y \( \forall b \in B \) existe
    un entorno \( U \) de \( b \) tal que \( p^{-1}(U) = \bigsqcup_{\alpha} V_\alpha \) tal que \( p|_{V_\alpha} \) es hómeo de \( V_\alpha \) con \( U \).
    Cada \( V_\alpha \) es una feta.
\end{definition}

\begin{ejercicio}{4}
    Dado \( p : E \rightarrow B \) un revestimiento. Mostrar que si \( B \) es compacto y \( p^{-1}(b) \) es finito para todo \( b \in B \) entonces \( E \) es compacto.
\end{ejercicio}
\begin{proof}
    Sea \( {\{ O_\lambda \}}_{\lambda \in \Lambda} \) un recubrimiento abierto cualquiera de \( E \). Queremos extraer un subrecubrimiento finito.

    Fijemos un punto \( b \in B \). Por hipótesis, la fibra \( p^{-1}(b) \) es finita, digamos \( p^{-1}(b) = \{ e_1, \ldots, e_n \} \).
    Para cada \( i \in \{1, \ldots, n\} \), como los \( O_\lambda \) cubren \( E \), existe algún \( O_{\lambda_i} \) tal que \( e_i \in O_{\lambda_i} \).

    Como \( p \) es un revestimiento, existe un entorno abierto \( U \subseteq B \) de \( b \) tal que \( p^{-1}(U) = \bigsqcup_{i=1}^{n} V_i \), donde cada \( V_i \) es un abierto de \( E \) con \( e_i \in V_i \), y \( p|_{V_i} : V_i \rightarrow U \) es un homeomorfismo.

    Para cada \( i \), consideremos la intersección \( W_i = V_i \cap O_{\lambda_i} \). Notemos que \( W_i \) es un entorno abierto de \( e_i \) en \( E \).
    Como \( p|_{V_i} \) es un homeomorfismo (y por lo tanto una función abierta), su imagen \( p(W_i) \) es un abierto en \( B \) que contiene a \( b \).

    Definimos \( N_b = \bigcap_{i=1}^{n} p(W_i) \). Como es una intersección finita de abiertos que contienen a \( b \), \( N_b \) es un entorno abierto de \( b \) en \( B \). Además, \( N_b \subseteq U \).

    Observemos qué ocurre al tomar preimagen de \( N_b \):
    \[ p^{-1}(N_b) = \bigsqcup_{i=1}^{n} {(p|_{V_i})}^{-1}(N_b) \]
    Por construcción, para cada \( i \), se tiene que \( {(p|_{V_i})}^{-1}(N_b) \subseteq {(p|_{V_i})}^{-1}(p(W_i)) = W_i \subseteq O_{\lambda_i} \).
    Es decir, toda la preimagen \( p^{-1}(N_b) \) está cubierta por la colección finita de abiertos \( \{ O_{\lambda_1}, \ldots, O_{\lambda_n} \} \).

    Repitiendo este proceso para cada \( b \in B \), obtenemos un recubrimiento abierto \( {\{ N_b \}}_{b \in B} \) del espacio base \( B \).
    Como \( B \) es compacto podemos extraer un subrecubrimiento finito \( \{ N_{b_1}, \ldots, N_{b_m} \} \).

    Finalmente, como estos abiertos cubren \( B \), sus preimágenes cubren \( E \):
    \[ E = p^{-1}(B) = p^{-1}\left( \bigcup_{j=1}^{m} N_{b_j} \right) = \bigcup_{j=1}^{m} p^{-1}(N_{b_j}) \]
    Cada \( p^{-1}(N_{b_j}) \) está cubierto por una cantidad finita de abiertos de nuestro recubrimiento original \( {\{ O_\lambda \}}_{\lambda \in \Lambda} \). Como hay una cantidad finita de estos conjuntos (\( m \) en total), concluimos que todo \( E \) se puede cubrir con una cantidad finita de abiertos de la familia original.

    Por lo tanto, \( E \) es compacto.
\end{proof}

\clearpage

\begin{ejercicio}{5}
    Dado \( p:E\rightarrow B \) un revestimiento con \( p(e_{0})=b_{0} \). Sean \( \gamma \) y \( \beta \) dos caminos en \( B \) tales que \( \gamma(1)=\beta(0) \).
    \begin{enumerate}
        \item Si \( \hat{\gamma} \) es la levantada de \( \gamma \) tal que \( \hat{\gamma}(0)=e_{0} \), \( \hat{\beta} \) es la levantada de \( \beta \) tal que \( \hat{\beta}(0)=\hat{\gamma}(1) \) y \( \widehat{\gamma*\beta} \) es la levantada de \( \gamma*\beta \) tal que \( \widehat{\gamma*\beta}(0)=e_{0} \) entonces \( \widehat{\gamma*\beta}=\hat{\gamma}*\hat{\beta} \).
        \item Si \( \hat{\gamma} \) es la levantada de \( \gamma \) tal que \( \hat{\gamma}(0)=e_{0} \) y \( \widehat{\gamma^{-1}} \) es la levantada de \( \gamma^{-1} \) tal que \( \widehat{\gamma^{-1}}(0)=\hat{\gamma}(1) \) entonces \( \hat{\gamma}*\widehat{\gamma^{-1}} \) es un \( e_{0} \)-rulo en \( E \).
    \end{enumerate}
\end{ejercicio}
\begin{proof}
    Sabemos que \( p(\hat{\gamma} \ast \hat{\beta}) = p(\hat{\gamma}) \ast p(\hat{\beta}) = \gamma \ast \beta \). Por la unicidad del levantamiento, pues comparten punto inicial, concluimos que \( \widehat{\gamma*\beta} = \hat{\gamma}*\hat{\beta} \).

    Para el segundo inciso, para que el camino \( \hat{\gamma}*\widehat{\gamma^{-1}} \) sea un \( e_0 \)-rulo en \( E \), debemos verificar que empieza y termina en \( e_0 \).
    El punto inicial es claro: \( (\hat{\gamma}*\widehat{\gamma^{-1}})(0) = \hat{\gamma}(0) = e_0 \).
    Para hallar el punto final, que es \( \widehat{\gamma^{-1}}(1) \), consideremos el camino \( \hat{\gamma} \) recorrido en sentido inverso, es decir, \( \hat{\gamma}^{-1}(t) = \hat{\gamma}(1-t) \).
    Notemos que \( p(\hat{\gamma}^{-1}(t)) = p(\hat{\gamma}(1-t)) = \gamma(1-t) = \gamma^{-1}(t) \), lo que significa que \( \hat{\gamma}^{-1} \) es una levantada de \( \gamma^{-1} \).
    Además, su punto inicial es \( \hat{\gamma}^{-1}(0) = \hat{\gamma}(1) \).
    Nuevamente, por la unicidad del levantamiento, la única levantada de \( \gamma^{-1} \) que empieza en \( \hat{\gamma}(1) \) es \( \widehat{\gamma^{-1}} \), por lo que \( \widehat{\gamma^{-1}} = \hat{\gamma}^{-1} \).
    En consecuencia, el punto final de nuestro camino es \( \widehat{\gamma^{-1}}(1) = \hat{\gamma}^{-1}(1) = \hat{\gamma}(0) = e_0 \).
    Como empieza y termina en \( e_0 \), concluimos que \( \hat{\gamma}*\widehat{\gamma^{-1}} \) es efectivamente un \( e_0 \)-rulo en \( E \).
\end{proof}

\vspace{.5cm}

\begin{ejercicio}{6}
    Sea \( p:E\rightarrow B \) un revestimiento con \( p(e_{0})=b_{0} \). Probar que:
    \begin{enumerate}
        \item El morfismo \( p_{*}:\pi_{1}(E,e_{0})\longrightarrow\pi_{1}(B,b_{0}) \) es un monomorfismo.
        \item Sea \( H \) la imagen por \( p_{*} \). Dado un \( b_{0} \)-rulo \( \gamma \) en \( B \), probar la equivalencia: \( [\gamma]\in H \iff \gamma \) se levanta a un rulo en \( E \).
        \item La \textit{correspondencia de levantamiento} induce una función inyectiva \[ \Phi:\frac{\pi_{1}(B,b_{0})}{H}\longrightarrow p^{-1}(b_{0}) \] que también será sobreyectiva en el caso que \( E \) sea arco-conexo.
        \item Si \( E \) es arco-conexo, puede definirse en \( p^{-1}(b_{0}) \) una estructura de grupo.
        \item Si \( E \) es simplemente conexo entonces el grupo \( p^{-1}(b_{0}) \) es isomorfo a \( \pi_{1}(B,b_{0}) \).
    \end{enumerate}
\end{ejercicio}
\begin{proof}
    Sea \( p \) el revestimiento y \( e_0 \in E \) tal que \( p(e_0) = b_0 \).
    \begin{enumerate}
        \item Queremos ver que \( p_* \) es un monomorfismo, veremos que \( \ker(p_*) = \{ [c_{e_0}] \} \).

              Dada \( [\gamma] \in \pi_1(E, e_0) \, : \, [\gamma] \in \ker(p_*) \). Entonces:
              \[ p_*({[\gamma]}) = [p \circ \gamma] = [c_{b_0}] \]
              donde \( c_{b_0} \) es el rulo constante en \( b_0 \).

              Que \( [p \circ \gamma] = [c_{b_0}] \) significa que existe una \( p \)-homotopía \( H : I \times I \rightarrow B \) entre el camino \( p \circ \gamma \) y el camino constante \( c_{b_0} \). Esta homotopía cumple las siguientes condiciones:
              \[
                  \begin{cases}
                      H(t,0) = (p \circ \gamma)(t) \\
                      H(t,1) = b_0                 \\
                      H(0,s) = b_0 \quad \text{y} \quad H(1,s) = b_0 \quad \forall s \in I
                  \end{cases}
              \]

              Como \( p \) es un revestimiento existe una única homotopía levantada \( \hat{H} : I \times I \rightarrow E \) tal que \( p \circ \hat{H} = H \) y cuyo camino inicial es \( \hat{H}(t,0) = \gamma(t) \).

              Estudiemos el comportamiento de \( \hat{H} \) en los extremos del intervalo:
              \begin{enumerate}
                  \item[1.] Para el extremo inicial \( t=0 \), el camino \( s \mapsto \hat{H}(0,s) \) se proyecta mediante \( p \) en el punto constante \( H(0,s) = b_0 \). Dado que las fibras de un revestimiento son espacios discretos y el intervalo \( I \) es conexo, la imagen de este camino debe ser un único punto. Como sabemos que empieza en \( \hat{H}(0,0) = \gamma(0) = e_0 \), concluimos que \( \hat{H}(0,s) = e_0 \) para todo \( s \in I \).
                  \item[2.] Aplicando exactamente el mismo argumento para el extremo final \( t=1 \), como \( H(1,s) = b_0 \) y \( \hat{H}(1,0) = \gamma(1) = e_0 \), se deduce que \( \hat{H}(1,s) = e_0 \) para todo \( s \in I \).
              \end{enumerate}
              Esto nos asegura que \( \hat{H} \) es una \( p \)-homotopía válida en \( E \) (mantiene los extremos fijos en \( e_0 \)).

              Finalmente, analicemos qué ocurre en el tiempo \( s=1 \). Por construcción, tenemos que \( p(\hat{H}(t,1)) = H(t,1) = b_0 \). Esto significa que el camino \( t \mapsto \hat{H}(t,1) \) es una levantada del camino constante \( c_{b_0} \) que comienza en \( \hat{H}(0,1) = e_0 \). Por la unicidad del levantamiento de caminos, la única levantada posible es el propio rulo constante en \( e_0 \), es decir, \( \hat{H}(t,1) = c_{e_0}(t) \).

              Por lo tanto, \( \hat{H} \) es una \( p \)-homotopía en \( E \) entre \( \gamma \) y \( c_{e_0} \), lo que demuestra que:
              \[ {[\gamma]} = [c_{e_0}] \]

              Como el único elemento en el núcleo es la clase de la identidad, concluimos que \( \ker(p_*) = \{ [c_{e_0}] \} \), lo que prueba que \( p_* \) es inyectiva y, por lo tanto, un monomorfismo.

        \item \( (\implies) \) Supongamos que \( [\gamma] \in H \). Por definición de la imagen, existe un \( e_0 \)-rulo \( f \) en \( E \) (es decir, \( f(0)=f(1)=e_0 \)) tal que \( p_*([f]) = [\gamma] \).

              Por la definición del morfismo inducido, esto significa que los caminos \( p \circ f \) y \( \gamma \) son p-homotópicos en \( B \):
              \[ {[p \circ f]} = [\gamma] \]

              Sea \( \hat{\gamma} \) la única levantada de \( \gamma \) que comienza en \( \hat{\gamma}(0) = e_0 \).
              Por otro lado, notemos que \( f \) es, por definición, la única levantada del camino \( p \circ f \) que comienza en \( e_0 \).

              Por el teorema de levantamiento de homotopías de caminos, como los caminos \( p \circ f \) y \( \gamma \) son homotópicos en la base y sus levantadas \( f \) y \( \hat{\gamma} \) comienzan en el mismo punto \( e_0 \), estas levantadas deben terminar exactamente en el mismo punto.

              Es decir:
              \[ \hat{\gamma}(1) = f(1) \]

              Pero como \( f \) es un \( e_0 \)-rulo, sabemos que \( f(1) = e_0 \). En consecuencia, \( \hat{\gamma}(1) = e_0 \).
              Como \( \hat{\gamma}(0) = \hat{\gamma}(1) = e_0 \), concluimos que \( \hat{\gamma} \) es un rulo en \( E \).

              \vspace{0.5cm}
              \( (\impliedby) \) Supongamos ahora que \( \gamma \) se levanta a un rulo \( \hat{\gamma} \) en \( E \) que comienza en \( e_0 \). Que sea un rulo significa que \( \hat{\gamma}(0) = \hat{\gamma}(1) = e_0 \), de modo que define un elemento válido para el grupo fundamental del espacio total, es decir, \( [\hat{\gamma}] \in \pi_1(E, e_0) \).

              Por la definición de levantada, sabemos que \( p \circ \hat{\gamma} = \gamma \). Si aplicamos el morfismo inducido \( p_* \) a la clase de \( \hat{\gamma} \), obtenemos:
              \[ p_*({[\hat{\gamma}]}) = [p \circ \hat{\gamma}] = [\gamma] \]
              Dado que \( [\gamma] \) es la imagen de un elemento de \( \pi_1(E, e_0) \) bajo \( p_* \), por definición de imagen concluimos que \( [\gamma] \in \text{Im}(p_*) = H \).

        \item TODO
        \item TODO
        \item TODO
    \end{enumerate}
\end{proof}

\vspace{1cm}

\begin{ejercicio}{7}
    Considerar \( p:\R \times \R^{+} \rightarrow \R^2 - \{(0,0)\} \) dado por \( p(t,r) = (r\cos(t), r\sin(t)) \).

    Asumiendo que \( p \) es un revestimiento, encontrar levantadas de las siguientes funciones: \( f(t)=(2-t,0) \) y \( g(t)=((1+t)\cos(t), (1+t)\sin(t)) \).
    Graficar las funciones y las levantadas.
\end{ejercicio}
\begin{proof}
    Buscamos funciones continuas \( \hat{f}(t) = (\theta(t), r(t)) \) y \( \hat{g}(t) = (\theta(t), r(t)) \) en \( E \) tales que \( p \circ \hat{f} = f \) y \( p \circ \hat{g} = g \). Supondremos el dominio estándar para caminos \( t \in [0,1] \).

    1. Levantadas de \( f(t) = (2-t, 0) \):
    Notemos que para \( t \in [0,1] \), el valor \( 2-t \) es estrictamente positivo (\( 1 \le 2-t \le 2 \)), por lo que el camino nunca toca el origen \( (0,0) \).

    Queremos encontrar \( \hat{f}(t) = (\theta(t), r(t)) \) tal que:
    \[ p(\theta(t), r(t)) = (r(t)\cos(\theta(t)), r(t)\sin(\theta(t))) = (2-t, 0) \]

    De la segunda componente del vector obtenemos:
    \[ r(t)\sin(\theta(t)) = 0 \]
    Como \( r(t) \in \mathbb{R}^{+} \), sabemos que \( r(t) > 0 \). Esto nos obliga a que \( \sin(\theta(t)) = 0 \), lo que significa que el ángulo debe ser un múltiplo entero de \( \pi \), es decir, \( \theta(t) = k\pi \) para algún \( k \in \mathbb{Z} \). Por continuidad de la función \( \theta \), este valor de \( k \) debe ser constante a lo largo de todo el camino.

    Sustituyendo esto en la primera componente:
    \[ r(t)\cos(k\pi) = 2-t \]
    Dado que tanto \( r(t) \) como \( 2-t \) son estrictamente mayores que cero para todo \( t \in [0,1] \), el signo de \( \cos(k\pi) \) debe ser positivo. Esto fuerza a que \( \cos(k\pi) = 1 \), lo cual solo ocurre si \( k \) es un número entero par. Definimos entonces \( k = 2m \) con \( m \in \mathbb{Z} \).

    Como \( \cos(2m\pi) = 1 \), despejamos el radio y obtenemos \( r(t) = 2-t \). Así, obtenemos una familia infinita de levantadas discretas indexadas por \( m \):
    \[ \hat{f}_m(t) = (2m\pi, 2-t), \quad \text{con } m \in \mathbb{Z} \]

    2. Levantadas de \( g(t) = ((1+t)\cos(t), (1+t)\sin(t)) \):
    Para \( t \in [0,1] \), la expresión \( 1+t \) también es estrictamente positiva (\( 1 \le 1+t \le 2 \)). Buscamos \( \hat{g}(t) = (\theta(t), r(t)) \) tal que:
    \[ (r(t)\cos(\theta(t)), r(t)\sin(\theta(t))) = ((1+t)\cos(t), (1+t)\sin(t)) \]

    Para hallar la componente del radio, podemos calcular la norma euclídea en ambos lados de la ecuación:
    \[ r(t) = \sqrt{{(1+t)}^2\cos^2(t) + {(1+t)}^2\sin^2(t)} = \sqrt{{(1+t)}^2} = |1+t| = 1+t \]

    Sustituyendo ahora \( r(t) = 1+t \) en nuestro sistema original y dividiendo ambos miembros por el factor común \( 1+t \), obtenemos de manera directa:
    \[ \cos(\theta(t)) = \cos(t) \quad \text{y} \quad \sin(\theta(t)) = \sin(t) \]
    La solución general para este sistema trigonométrico nos indica que los ángulos deben diferir en múltiplos enteros de una vuelta completa. Por lo tanto, \( \theta(t) = t + 2k\pi \) para algún \( k \in \mathbb{Z} \) fijo por continuidad.

    De este modo, obtenemos la familia infinita de levantadas para \( g(t) \):
    \[ \hat{g}_k(t) = (t + 2k\pi, 1+t), \quad \text{con } k \in \mathbb{Z} \]

    Para realizar los gráficos requeridos, la descripción del comportamiento de las curvas es la siguiente:
    \begin{itemize}
        \item \textbf{En el espacio base \( B = \R^2 - \{(0,0)\} \):}
              El camino \( f(t) \) es simplemente un segmento de recta horizontal situado sobre el eje \( x \) positivo que se recorre de derecha a izquierda, comenzando en el punto \( (2,0) \) y terminando en \( (1,0) \).
              El camino \( g(t) \) representa un arco de una espiral de Arquímedes. Comienza en el punto \( (1,0) \) cuando \( t=0 \) y va girando en sentido antihorario a la vez que se aleja del origen, finalizando en el punto \( (2\cos(1), 2\sin(1)) \) cuando \( t=1 \).

        \item \textbf{En el espacio total \( E = \R \times \R^{+} \) (visto como un plano cartesiano donde el eje horizontal es \( t \) y el eje vertical es \( r \)):}
              Las levantadas \( \hat{f}_m(t) \) son infinitos segmentos de recta completamente verticales localizados sobre las abscisas \( t = 2m\pi \). A medida que \( t \) avanza de \( 0 \) a \( 1 \), estas líneas descienden verticalmente desde la altura \( r=2 \) hasta la altura \( r=1 \).
              Las levantadas \( \hat{g}_k(t) \) son infinitos segmentos de recta oblicuos que tienen una pendiente exacta de 45 grados (ya que tanto la coordenada angular como el radio aumentan linealmente a la misma velocidad con respecto a \( t \)). Cada una de estas líneas rectas comienza en el punto \( (2k\pi, 1) \) y avanza hacia la derecha y hacia arriba hasta alcanzar el punto \( (1+2k\pi, 2) \).
    \end{itemize}
\end{proof}

\vspace{0.5cm}

\begin{ejercicio}{8}
    Probar que dos espacios homeomorfos son homotópicamente equivalentes; pero no vale la recíproca.
\end{ejercicio}
\begin{proof}
    Que dos espacios \( X \) e \( Y \) sean homeomorfos significa que existe un homeomorfismo \( f : X \to Y \) tal que \( f \) es continua, biyectiva y su inversa \( f^{-1} \) también es continua.
    Para demostrar que son homotópicamente equivalentes, es decir, que existen funciones continuas \( g : X \to Y \) e \( h : Y \to X \)
    tales que \( h \circ g \simeq id_X \) y \( g \circ h \simeq id_Y \), podemos tomar \( g = f \) y \( h = f^{-1} \).
    Entonces:
    \[
        h \circ g = f^{-1} \circ f = id_X \quad \text{y} \quad g \circ h = f \circ f^{-1} = id_Y
    \]
    Como \( id_X \) y \( id_Y \) son triviales, se concluye que \( X \) e \( Y \) son homotópicamente equivalentes.

    La recíproca no vale pues \( \R^n \) es contractil y, por lo tanto, homotópicamente equivalente a un punto, pero no es homeomorfo a un punto.
\end{proof}

\vspace{0.5cm}

\begin{definition}
    Sea \( A \subseteq X \) un subespacio. Decimos que \( A \) es un retracto por deformación de \( X \) si existe una homotopía \( H : X \times [0,1] \to X \) tal que:
    \begin{enumerate}
        \item \( H(x,0) = id_X(x) = x \) para todo \( x \in X \).
        \item \( H(x,1) \in A \) para todo \( x \in X \).
        \item \( H(a,t) = a \) para todo \( a \in A \) y todo \( t \in [0,1] \).
    \end{enumerate}
\end{definition}

\begin{ejercicio}{9}
    El ecuador de la esfera \( S^{2} \) es un retracto por deformación de la esfera sin sus 2 polos.
\end{ejercicio}
\begin{proof}
    Sea \( X = S^2 \setminus \{N, S\} \) la esfera sin los polos, donde \( N = (0,0,1) \) y \( S = (0,0,-1) \).
    El ecuador es el subespacio \( E = \{ (x,y,0) \in S^2 \} \).

    La idea geométrica es empujar la coordenada \( z \) hacia cero progresivamente, al mismo tiempo que normalizamos el vector resultante para que la imagen permanezca en la esfera. Definimos la homotopía como:
    \[
        H((x,y,z), t) = \frac{ (x, y, (1-t)z) }{ \left \| (x, y, (1-t)z) \right \| }
    \]

    Verifiquemos que \( H \) está bien definida y cumple las propiedades:
    \begin{itemize}
        \item \textbf{Continuidad y buena definición:} Como el punto de partida \( (x,y,z) \) pertenece a \( X \), sabemos que no es ninguno de los polos, lo que implica que \( x^2 + y^2 > 0 \). Por lo tanto, para cualquier \( t \in [0,1] \), el denominador cumple que \( \sqrt{x^2+y^2+(1-t)^2z^2} \ge \sqrt{x^2+y^2} > 0 \), y nunca se anula. Como la norma de \( H((x,y,z),t) \) es exactamente 1 por construcción, la imagen cae en \( S^2 \). Finalmente, como las primeras dos coordenadas no pueden ser ambas simultáneamente cero, la imagen nunca puede ser \( (0,0,\pm 1) \), asegurando que cae en \( X \).

        \item \textbf{Condición inicial (\( t=0 \)):}
              \[ H((x,y,z), 0) = \frac{(x,y,z)}{\sqrt{x^2+y^2+z^2}} = (x,y,z) \]
              Esto se debe a que, como \( (x,y,z) \) está en la esfera, \( x^2+y^2+z^2 = 1 \).

        \item \textbf{Condición final (\( t=1 \)):}
              \[ H((x,y,z), 1) = \left( \frac{x}{\sqrt{x^2+y^2}}, \frac{y}{\sqrt{x^2+y^2}}, 0 \right) \]
              La coordenada \( z \) se hace cero, por lo que la imagen claramente pertenece al ecuador \( E \).

        \item \textbf{Puntos fijos en el ecuador (\( z=0 \)):} Si tomamos un punto que ya está en el ecuador, su coordenada \( z \) es nula, y además \( x^2+y^2=1 \). Evaluando para cualquier tiempo \( t \):
              \[ H((x,y,0), t) = \left( \frac{x}{\sqrt{x^2+y^2}}, \frac{y}{\sqrt{x^2+y^2}}, 0 \right) = (x,y,0) \]
    \end{itemize}
    Por lo tanto, \( H \) es un retracto por deformación de \( X \) en \( E \).
\end{proof}

\vspace{0.5cm}

\begin{ejercicio}{10}
    Dadas \( f \) y \( g \in C(X,S^{2}) \). Probar que si \( |f(x)-g(x)|<2 \) para todo \( x\in X \) entonces \( f \) y \( g \) son homotópicas.
\end{ejercicio}
\begin{proof}
    El hecho de que \( |f(x) - g(x)| < 2 \) para todo \( x \in X \) implica que \( f(x) \) y \( g(x) \) no son antipodales en la esfera \( S^2 \).
    Esto nos permite definir una homotopía lineal entre \( f \) y \( g \) a lo largo del segmento de línea que conecta \( f(x) \) y \( g(x) \) en el espacio euclidiano \( \R^3 \).
    Defino: \( F : X \times [0, 1] \to S^2 \) por:
    \[
        F(x, t) = \frac{(1-t)f(x) + t g(x)}{\| (1-t)f(x) + t g(x) \| }
    \]
    Donde \( F(x, 0) = f(x) \) y \( F(x, 1) = g(x) \).
    La condición \( |f(x) - g(x)| < 2 \) asegura que el vector \( (1-t)f(x) + t g(x) \) nunca es cero, por lo que la normalización es válida y \( F \) es continua.
    Por lo tanto, \( F \) es una homotopía entre \( f \) y \( g \), demostrando que son homotópicas.
\end{proof}

\vspace{0.5cm}

\begin{ejercicio}{11}
    Si \( X \) es un RD de \( Y \) e \( Y \) es un RD de \( Z \) entonces \( X \) es un RD de \( Z \).
\end{ejercicio}
\begin{proof}
    Notemos que como \( X \) es RD de \( Y \) existe una homotopía \( H_1 : Y \times [0, 1] \to Y \) tal que \( H_1(y,0) = y \), \( H_1(y,1) \in X \) y \( H_1(x,t) = x \) para todo \( x \in X \).
    De manera análoga, como \( Y \) es RD de \( Z \), existe una homotopía \( H_2 : Z \times [0, 1] \to Z \) tal que \( H_2(z,0) = z \), \( H_2(z,1) \in Y \) y \( H_2(y,t) = y \) para todo \( y \in Y \).

    Consideremos la homotopía compuesta \( H : Z \times [0, 1] \to Z \) definida por:
    \[
        H(z, t) = \begin{cases}
            H_2(z, 2t)             & \text{si } 0 \le t \le \frac{1}{2} \\
            H_1(H_2(z, 1), 2t - 1) & \text{si } \frac{1}{2} < t \le 1
        \end{cases}
    \]
    Verifiquemos que \( H \) cumple las condiciones de un retracto por deformación de \( Z \) en \( X \):
    \begin{itemize}
        \item Para \( t=0 \), \( H(z, 0) = H_2(z, 0) = z \), cumpliendo la condición inicial.
        \item Para \( t=1 \), \( H(z, 1) = H_1(H_2(z, 1), 1) \). Como \( H_2(z, 1) \in Y \), aplicando \( H_1 \) nos lleva a \( H_1(H_2(z, 1), 1) \in X \), cumpliendo la condición final.
        \item Para \( x \in X \), notemos que \( H_2(x, t) = x \) para todo \( t \), ya que \( x \in X \subseteq Y \). Por lo tanto, \( H(x, t) = H_1(x, 2t - 1) = x \) para \( t > \frac{1}{2} \), y para \( t \le \frac{1}{2} \), \( H(x, t) = H_2(x, 2t) = x \). En ambos casos, \( H(x, t) = x \), cumpliendo la condición de puntos fijos en \( X \).
    \end{itemize}
    Por lo tanto, \( H \) es una homotopía que demuestra que \( X \) es un retracto por deformación de \( Z \).
\end{proof}

\vspace{0.5cm}

\begin{ejercicio}{12}
    El disco unidad es un RD de \(\R^n\).
\end{ejercicio}
\begin{proof}
    Como venimos haciendo, consideremos la siguiente homotopía y veamos que cumple la definición de retracto por deformación:
    \[
        H : \R^n \times [0, 1] \to \R^n, \quad H(x, t) = (1-t)x + t\frac{x}{\|x\|}
    \]
    Verifiquemos que \( H \) cumple las condiciones de un retracto por deformación:
    \begin{itemize}
        \item Para \( t=0 \), \( H(x, 0) = x \), cumpliendo la condición inicial.
        \item Para \( t=1 \), \( H(x, 1) = \frac{x}{\|x\|} \), que pertenece al disco unidad, cumpliendo la condición final.
        \item Para \( x \) en el disco unidad \( \| x \| = 1 \implies H(x, t) = x \) para todo \( t \), cumpliendo la condición de puntos fijos.
    \end{itemize}
    Por lo tanto, \( H \) es una homotopía que demuestra que el disco unidad es un retracto por deformación de \(\R^n\).
\end{proof}

\vspace{0.5cm}

\begin{ejercicio}{13}
    Si \( Y \) y \( Y^{\prime} \) son RD de \( X \) y \( X^{\prime} \) respectivamente, entonces \( Y\times Y^{\prime} \) es un RD de \( X\times X^{\prime} \).
\end{ejercicio}
\begin{proof}
    Sea \( H : X \times [0, 1] \to X \) una homotopía que demuestra que \( Y \) es un retracto por deformación de \( X \), y sea \( H' : X' \times [0, 1] \to X' \) una homotopía que demuestra que \( Y' \) es un retracto por deformación de \( X' \).

    Definimos la homotopía \( K : (X \times X') \times [0, 1] \to X \times X' \) por:
    \[
        K((x, x'), t) = (H(x, t), H'(x', t))
    \]

    Verifiquemos que \( K \) cumple las condiciones de un retracto por deformación:
    \begin{itemize}
        \item Es continua pues lo es coordenada a coordenada.
        \item Para \( t=0 \), \( K((x, x'), 0) = (H(x, 0), H'(x', 0)) = (x, x') \), cumpliendo la condición inicial.
        \item Para \( t=1 \), \( K((x, x'), 1) = (H(x, 1), H'(x', 1)) \in Y \times Y' \), cumpliendo la condición final.
        \item Para \( (y, y') \in Y \times Y' \), tenemos:
              \[
                  K((y, y'), t) = (H(y, t), H'(y', t)) = (y, y')
              \]
              para todo \( t \in [0, 1] \), cumpliendo la condición de puntos fijos.
    \end{itemize}

    Por lo tanto, \( K \) es una homotopía que demuestra que \( Y\times Y^{\prime} \) es un retracto por deformación de \( X\times X^{\prime} \).
\end{proof}

\vspace{0.5cm}

\begin{ejercicio}{14}
    Probar que si \( f:S^{n}\rightarrow S^{n} \) es continua y no tiene puntos fijos entonces \( f \) es homotópica a la simetría \( x\rightarrow-x \).
\end{ejercicio}
\begin{proof}
    TODO
\end{proof}

De la práctica de 2025.

\begin{ejercicio}{15}
    Sean \(f, g \in C(X, \C \setminus \{ 0 \}) \) tales que \( |f(x) - g(x)| < |f(x)| \quad \forall x \in X \implies f \) homotópica a \( g \).
\end{ejercicio}
\begin{proof}
    Consideremos la homotopía \( H : X \times [0, 1] \to \C \setminus \{ 0 \} \) dada por:
    \[
        H(x, t) = (1-t)f(x) + t g(x)
    \]
    Sólo debemos chequear que \( H(x, t) \neq 0 \) para todo \( x \in X \) y \( t \in [0, 1] \).
    En efecto, como \begin{align*}
        |H(x, t)| & = |(1-t)f(x)+t g(x)|        \\
                  & = |f(x) + t(g(x) - f(x))|   \\
                  & \ge |f(x)| - t|g(x) - f(x)| \\
                  & = |f(x)| - t|f(x) - g(x)|   \\
    \end{align*}
    Como \( |f(x) - g(x)| < |f(x)| \) y \( t \in [0, 1] \), se deduce que:
    \[
        |H(x, t)| > |f(x)| - |f(x) - g(x)| > 0
    \]
    Por lo tanto, \( H(x, t) \in \C \setminus \{ 0 \} \) para todo \( x \in X \) y \( t \in [0, 1] \), lo que demuestra que \( f \) es homotópica a \( g \).
\end{proof}

\end{document}